\documentclass[11pt]{article}
\usepackage{theme}
\usepackage{shortcuts}
\usepackage[backend=biber, sorting=none]{biblatex}
\usepackage{hyperref}
\usepackage[title]{appendix}
\usepackage{subcaption}
\usepackage{float}
\usepackage{siunitx}
\usepackage{array}
\usepackage{tabularx}   % For auto-adjusting column widths
\usepackage[figurename=Fig.,font=small,labelfont=bf,labelsep=space]{caption}
\renewcommand{\figureautorefname}{Fig.}
\addbibresource{Library.bib}
\def\mathdefault#1{#1}
\setlist{leftmargin=5mm}
\setlength\itemsep{0mm}
% Document parameters
% Document title
\title{Benchmarking Foundation Models using Satellite Image Time
Series for Land Monitoring - MVA 2025/2026}
\author{
  Jean-François Guerrero
  \email{jean-francois.guerrero@ens-paris-saclay.fr} \\ % student 1
  Simon Sassi \email{sassisimon84@gmail.com} % student 2
}

\begin{document}
\maketitle
\newpage
\section{Introduction}
\section{Dataset}
We use the dataset PASTIS (Panoptic Agricultural Satellite TIme Series)
There are 2433 sequences of multi-spectral of shape 10*128*128.Each
sequences contains between 38 and 61 images. The satellite THEIA
provides observations.
The 10 channels correspond to the non-atmospheric spectral
bands of the Sentinel-2 satellite, after atmospheric correction and
re-sampling at a spatial resolution of 10 meters per pixel.
Observatios are taken in several region of France with different
climates, (Normandie, Bourgogne,Alsace, Provence). Almost 30 $\%$ of
pictures are covered by cloud. \cite{dumeur}

Each pixel is associated with a semantic label of 18 crop labels
% TODO : problème d'imbalance des classes MIOU, mettre que les
% classes ne sont pas balancées, pas très fair,
% TODO : F1 par classe sur le scarce ?
% TODO : Nombre de paramètres par modèles

\section{Methods}
\section{Results}
\section{Summary of findings}

\newpage
\printbibliography

\newpage
\begin{appendices}
  \section{Method}
  % \begin{figure}
  %   \centering
  %   \includegraphics[width=\textwidth]{figures/fair_attri_architecture_schema.pdf}
  %   \caption{Overview of our training framework.}
  %   \label{fig:fair_attri_architecture_schema}
  % \end{figure}
  \section{Data}

  \section{Tables}
  \begin{table}[ht] % Ajout d'un spécificateur de position
    \centering % Pour centrer le tableau sur la page
    \caption{Test metrics}
    \label{tab:test_metrics}
    \begin{tabular}{|l| ccc|}
      \hline
      Model & mIoU & F1 Macro & Overall accuracy \\ % Ajout d'un
      % titre de colonne
      \hline
      AlphaEarth & 0.5287 & 0.6526 & 0.8417 \\
      Tessera    & 0.5607 & 0.7018 & 0.8257 \\
      ALISE      & 0.6009 & \textbf{0.7235} & 0.8651 \\
      U-TAE &  \textbf {0.7219} & - &  \textbf{0.9170} \\
      \hline
    \end{tabular}
  \end{table}

  \begin{table}[ht]
    \centering
    % Réduit légèrement la taille pour que les noms longs rentrent
    \caption{Test metrics with scarce data (30 patches per fold)}
    \label{tab:test_metrics_scarce}
    \begin{tabular}{|l |ccc|}
      \hline
      Model & mIoU & F1 Macro & Overall accuracy \\
      \hline
      AlphaEarth & 0.2867 & 0.3755 & 0.7267 \\
      Tessera & 0.4387 & \textbf{0.5765} & 0.7621 \\
      ALISE & \textbf{0.4482} & 0.5729 & \textbf{0.7898} \\
      U-TAE & 0.3869 & 0.5115 & 0.7043 \\
      \hline
    \end{tabular}
  \end{table}
\end{appendices}
\end{document}
